\chapter{Finite cyclic reductions}
\label{mht:ch:finite-reductions}

A finite collection of curve modes can retain the kinetic pairing
without retaining the product. The projected product then has an
associator, and the eliminated modes determine the higher operations
that cancel it. We construct this cancellation on the smooth Neumann
complex. A six-dimensional example already has nonzero ternary and
quartic operations. We then vary the residual matrix and transport the
finite subspace so that its pairing stays nondegenerate through every
eigenvalue collision.

\section{Finite subspaces and the polar trace}

Retain $E$, $d_E$, $\tau_E$, $p_0$, $p_1$, and $h$ from
Chapter~\ref{mht:ch:polar-neumann}. Throughout this chapter choose
\[
 f_0=\frac{z}{1+|z|^2},\qquad \eta_0=d_Ef_0,
\]
and keep the circular prescription for the double-pole trace.
There are two elementary restrictions on a finite approximation.

\begin{proposition}\label{mht:prop:finite-restrictions}
On a connected smooth manifold, every finite-dimensional unital
subalgebra of complex-valued smooth functions consists of constants.
On $C=\mathbb P^1$, no operator
$R:\Omega^{0,1}(C)\to\Omega^{0,1}(C)$ with a globally smooth
Schwartz kernel preserves $\tau_E$.
\end{proposition}

\begin{proof}
The powers of a function in a finite-dimensional algebra satisfy a
nonzero polynomial relation. The function therefore takes values in
a finite set. Connectedness makes it constant.

The polar trace is a distribution of finite order. Applying it in
the output variable of a smooth kernel gives a smooth covector in
the input variable. If $\tau_ER=\tau_E$, that covector would equal
$\kappa z^{-2}dz/(2\pi i)$ away from $o$. This coefficient has no
smooth extension across $o$.
\end{proof}

Thus a finite space containing a nonconstant scalar mode cannot
retain the pointwise product by restriction. A projection preserving
the trace must also retain its distributional input covector. Neither
restriction prevents a cyclic contraction.

The acyclic part of the complex is
\[
 V=\ker p_0\subset E^0,\qquad W=\ker p_1\subset E^1,
 \qquad d_E:V\xrightarrow{\sim}W,
\]
with inverse $h|_W$. On $V$ define
\begin{equation}\label{mht:eq:acyclic-symplectic-form}
 \omega_0(v,w)=\tau_E(vd_Ew).
\end{equation}

\begin{lemma}\label{mht:lem:finite-symplectic}
The form $\omega_0$ is continuous, skew and weakly nondegenerate.
Let $\mathcal R$ be the restrictions to the round sphere of complex
polynomials in its three Cartesian coordinates. Then
\[
 E_{\mathrm{alg}}^0=\mathcal R\cap\ker j,
 \qquad E_{\mathrm{alg}}^1=d_E\mathcal R
\]
form a dense subalgebra and subcomplex of $E$, preserved by $h$ and
the residual projections. The restriction of $\omega_0$ to
$V_{\mathrm{alg}}=V\cap\mathcal R$ is nondegenerate.
\end{lemma}

\begin{proof}
For $v,w\in E^0$, polar Stokes applied to $vw$ gives
$\tau_E(vd_Ew)=-\tau_E((d_Ev)w)$. If $v\in V$ annihilates
$V$ under this form, it pairs to zero with $W=d_EV$. The equality
$p_0(v)=0$ also gives $\tau_E(v\eta_0)=0$. Hence
$\tau_E(v\eta)=0$ for every smooth one-form $\eta$. Tests supported
away from $o$ imply $v=0$ there, and smoothness gives $v=0$ on $C$.

The space $\mathcal R$ consists of finite sums of spherical
harmonics. Smooth functions and forms have convergent expansions in
rotation types in every smooth seminorm. The Dolbeault differential
and the round Dolbeault Laplacian preserve these types. Since
$H^{0,1}(C)=0$, the Dolbeault inverse sends a finite sum of form
types to a finite sum of function types. Subtracting its value at
$o$ gives $G$ and preserves this property.

Products of finite sums
of types are again finite sums, also for a function times a form.
This proves closure under products and $G$, hence under $h$ and $p$.

For a sequence $f_n\in\mathcal R$ converging to a Neumann function,
replace $f_n$ by $f_n-j(f_n)f_0$. Subtracting its $p_0$ projection
also imposes the condition defining $V$. These continuous corrections
prove the required density. A vector of $V_{\mathrm{alg}}$ that
annihilates $V_{\mathrm{alg}}$ annihilates its closure $V$, so the
first paragraph proves nondegeneracy on $V_{\mathrm{alg}}$.
\end{proof}

Choose a basis of the countable-dimensional space $V_{\mathrm{alg}}$.
Starting with a finite nondegenerate subspace, take the first basis
vector outside it and subtract its symplectic projection. The nonzero
remainder has a partner by Lemma~\ref{mht:lem:finite-symplectic}.
Subtract the partner's projection as well and adjoin the resulting
two-dimensional symplectic space. Iteration gives
\[
 V_1\subset V_2\subset\cdots,
 \qquad \dim V_K=2K,
 \qquad \bigcup_KV_K=V_{\mathrm{alg}}.
\]

For a basis $(v_a)_{a=1}^{2K}$ of $V_K$, its symplectic projection is
\begin{equation}\label{mht:eq:finite-symplectic-projection}
 \Omega_{ab}=\omega_0(v_a,v_b),\qquad
 \Pi_Kv=\sum_{a,b}v_a(\Omega^{-1})_{ab}\omega_0(v_b,v).
\end{equation}
Only a finite matrix is inverted.

\begin{definition}\label{mht:def:finite-neumann-complex}
The finite Neumann complex and its projection are
\begin{align}
 E_K^0&=\mathbb C1\oplus V_K,&
 E_K^1&=\mathbb C\eta_0\oplus d_EV_K,\notag\\
 \mathsf p_K^0f&=p_0(f)1+\Pi_K(f-p_0(f)1),&
 \mathsf p_K^1\eta&=p_1(\eta)\eta_0+d_E\Pi_Kh\eta.
 \label{mht:eq:finite-neumann-projection}
\end{align}
Write $i_K:E_K\hookrightarrow E$ for inclusion and set
\[
 H_K\eta=(1-\Pi_K)h\eta,\qquad H_K(E^0)=0.
\]
\end{definition}

\begin{theorem}\label{mht:thm:finite-cyclic-contraction}
The pairing $\tau_E(xy)$ is nondegenerate on $E_K$. The maps above
give a continuous cyclic strong deformation retract:
\[
 \mathsf p_Ki_K=1,\qquad d_EH_K+H_Kd_E=1-i_K\mathsf p_K,
 \qquad H_K^2=H_Ki_K=\mathsf p_KH_K=0.
\]
Projection is adjoint to inclusion, and both preserve the unit and
trace. For $K\leq L$, the induced contraction from $E_L$ to $E_K$
has homotopy
\begin{equation}\label{mht:eq:finite-nested-homotopy}
 H_{LK}=(\Pi_L-\Pi_K)h|_{E_L},\qquad
 H_K=H_L+i_LH_{LK}\mathsf p_L.
\end{equation}
\end{theorem}

\begin{proof}
The decomposition
\[
 E=(\mathbb C1\oplus\mathbb C\eta_0)
      \oplus(V\xrightarrow{d_E}W)
\]
is orthogonal for the cyclic pairing. The first summand has
$\langle1,\eta_0\rangle=\kappa$. Identifying $W$ with $V$ by $h$
identifies the pairing between $V_K$ and $d_EV_K$ with $\Omega$.
This proves nondegeneracy on $E_K$.

The symplectic projection $\Pi_K$ is self-adjoint for $\omega_0$.
Its transport to $W$ is $d_E\Pi_Kh$, and the inverse differential
on the orthogonal complement is $(1-\Pi_K)h$. This proves the
homotopy identity and its side conditions. The same orthogonality
and polar Stokes prove cyclicity of $H_K$. The unit and its dual
$\eta_0$ are retained, giving preservation of the trace. All the
coefficient functionals in the projection are continuous.

Nested symplectic projections satisfy
$\Pi_K\Pi_L=\Pi_L\Pi_K=\Pi_K$. Decompose the acyclic summand
into $V_K$, its complement in $V_L$, and the complement of $V_L$.
The inverse differential on these summands gives
\eqref{mht:eq:finite-nested-homotopy}.
\end{proof}

\section{A six-dimensional example}

The first finite complex can be made explicit. Put $r=|z|^2$ and
\begin{equation}\label{mht:eq:finite-two-modes}
 v=\frac r{1+r}-\frac12,
 \qquad w=\frac{zr}{(1+r)^2}+\frac{\bar z}{1+r},
 \qquad V_1=\operatorname{span}\{v,w\}.
\end{equation}

Both functions are smooth in the two sphere charts, satisfy $j=0$,
and have zero $p_0$ projection. Their pairings are
\begin{equation}\label{mht:eq:finite-two-pairings}
 \Omega=\kappa\begin{pmatrix}0&1/6\\-1/6&0\end{pmatrix},
 \qquad
 S=\bigl(\tau_E(\eta_0v_av_b)\bigr)
       =\kappa\begin{pmatrix}1/12&0\\0&1/6\end{pmatrix}.
\end{equation}
These values follow from one radial integral. For integers
$0\leq a\leq k$ and $0\leq b\leq k-2$, the indicated form is smooth
at infinity, and angular integration gives
\begin{equation}\label{mht:eq:finite-radial-moment}
 \kappa^{-1}\tau_E\left(
      \frac{z^a\bar z^b}{(1+|z|^2)^k}d\bar z\right)
 =\begin{cases}
 -\dfrac{b!(k-b-2)!}{(k-1)!},&a-b=2,\\
 0,&a-b\ne2.
 \end{cases}
\end{equation}
In the nonzero case the remaining integral is
$-\int_0^\infty r^b(1+r)^{-k}\,dr$.

Write $\pi_1$ for the projection to the complex $E_1$ in both
degrees, and put $m_2(x,y)=\pi_1(xy)$. Substituting
\eqref{mht:eq:finite-two-pairings} into
\eqref{mht:eq:finite-neumann-projection} gives
\begin{equation}\label{mht:eq:finite-product-table}
 m_2(v,v)=\frac1{12},\qquad
 m_2(w,w)=\frac16,\qquad m_2(v,w)=0.
\end{equation}
Consequently
\begin{equation}\label{mht:eq:finite-associator}
 m_2(m_2(v,v),w)-m_2(v,m_2(v,w))=\frac w{12}.
\end{equation}
For the projected matrix commutator, the cyclic Jacobi sum
$[[x,y],z]+[[y,z],x]+[[z,x],y]$ on
$x=vE_{12}$, $y=vE_{21}$, $z=wE_{12}$ is $wE_{12}/6$.
Thus keeping the projected binary operation alone loses the gauge
identity already on these smooth modes.

\section{Higher operations from the eliminated modes}

Let $\sigma:E\to E[1]$ have degree $-1$. On the shifted complex use
\[
 b_1=-\sigma d_E\sigma^{-1},\qquad
 b_2(\sigma x,\sigma y)=(-1)^{|x|}\sigma(xy),\qquad
 \mathbf H_K=-\sigma H_K\sigma^{-1}.
\]

Shift $i_K$ and $\mathsf p_K$ as well. Define maps on shifted tensor powers
by the recursion
\begin{align}
 I_1&=i_K,&\beta_1&=b_1|_{E_K[1]},\notag\\
 \lambda_n&=\sum_{r=1}^{n-1}b_2(I_r\otimes I_{n-r}),&
 I_n&=-\mathbf H_K\lambda_n,&
 \beta_n&=\mathsf p_K\lambda_n\quad(n\geq2).
 \label{mht:eq:finite-tree-recursion}
\end{align}
Tensor products of maps have their Koszul signs. Each $I_n$ has
degree zero on the shifted spaces, and each $\beta_n$ has degree
one. At a fixed arity the recursion is a finite sum over planar
binary trees.

\begin{theorem}\label{mht:thm:finite-cyclic-transfer}
The maps $\beta_n$ extend to a square-zero coderivation on the
conilpotent tensor coalgebra $T^c(E_K[1])$. They define a strictly
unital cyclic $A_\infty$ algebra, and the maps $I_n$ give an
$A_\infty$ quasi-isomorphism to $E$. Tensoring with matrices and
antisymmetrizing gives cyclic $L_\infty$ operations with the full
matrix trace.

For the nested contractions of
\eqref{mht:eq:finite-nested-homotopy}, successive transfer through
$E_L$ gives the same operations on $E_K$ as direct transfer.
The tree comparisons are compatible with every further nested
elimination.
\end{theorem}

\begin{proof}
The differential, Leibniz rule and associativity are the coderivation
identities
\[
 b_1^2=0,\qquad b_1b_2+b_2b_1=0,\qquad b_2^2=0.
\]
The shifted homotopy satisfies
$b_1\mathbf H_K+\mathbf H_Kb_1=1-i_K\mathsf p_K$.

Move $b_1$ through a tree in
\eqref{mht:eq:finite-tree-recursion}. Its terms at vertices cancel
by the Leibniz rule. At an internal edge, replace the commutator
with $\mathbf H_K$ by $1-i_K\mathsf p_K$. The identity contracts the edge.
The two resolutions of the resulting ternary vertex cancel by
associativity. The $i_K\mathsf p_K$ term cuts the tree into two transferred
trees.

Summing all cuts gives $\beta^2=0$ at the root and
$(b_1+b_2)I=I\beta$ at the inclusion. The linear component is a
quasi-isomorphism by Theorem~\ref{mht:thm:finite-cyclic-contraction}.

Pair the root with one additional input. Adjointness of $\mathsf p_K$ and
cyclicity of $H_K$ allow the root to move across every edge, with
the graded sign of the pairing. The resulting unrooted tensor is
cyclic.

This is the cyclic tree transfer associated with an abstract
Hodge decomposition.\footnote{J.~Chuang and A.~Lazarev,
\emph{Abstract Hodge decomposition and minimal models for cyclic
algebras}, \href{https://arxiv.org/abs/0810.2393}{arXiv:0810.2393}.
See Definition~2.1 and Theorem~4.2.
Here the retained differential need not vanish.}

Suppose a unit enters a tree of arity at least three. Its neighboring
vertex sends either a retained input into $H_Ki_K=0$ or a homotopy
output into $H_K^2=0$. If that vertex is the root, the
same term vanishes by $\mathsf p_KH_K=0$. The binary operation preserves
the unit. Antisymmetrization of the associative identities gives
the Lie identities and preserves the cyclic tensor.

For successive transfer, expand each internal homotopy using
\eqref{mht:eq:finite-nested-homotopy}. Mark the edge by its chosen
summand. The maximal subtrees whose edges use $H_L$ are exactly
the transferred vertices on $E_L$. The remaining edges use $H_{LK}$.
This is a bijection of signed planar trees with those of successive
transfer. With three or more cutoffs, grouping the connected subtrees
in either order gives the same marked tree.
\end{proof}

The missing operations in the six-dimensional example can now be
calculated. In the next formulas, $\beta_n(x_1,\ldots,x_n)$ denotes
$\sigma^{-1}\beta_n(\sigma x_1,\ldots,\sigma x_n)$, with this
convention fixed rather than an additional desuspension sign.

For inputs of degrees $a,b,c$, the cubic recursion reads
\[
 \beta_3(x,y,z)=(-1)^{a+b}\pi_1\bigl(xH_1(yz)\bigr)
                      -(-1)^b\pi_1\bigl(H_1(xy)z\bigr).
\]
It gives
\begin{align}
 \beta_3(v,v,d_Ev)&=v/30,&
 \beta_3(v,v,d_Ew)&=w/30,\notag\\
 \beta_3(d_Ev,v,w)&=w/60,&
 \beta_3(v,d_Ev,w)&=-w/30.
 \label{mht:eq:finite-cubic-products}
\end{align}
For example, direct Green integration and
\eqref{mht:eq:finite-radial-moment} give
\[
 H_1(vd_Ev)=\frac{r^2-4r+1}{12(1+r)^2},\qquad
 \pi_1\bigl(vH_1(vd_Ev)\bigr)=v/30.
\]
On $(v,v,w)$ the complete cubic identity reduces to
\[
 \frac1{12}-\frac1{60}-\frac1{30}-\frac1{30}=0.
\]
The unary differential is essential here: $v$ and $w$ are not
closed cohomology classes.

The fourth operation also survives:
\begin{equation}\label{mht:eq:finite-quartic-products}
 \begin{aligned}
 \beta_4(d_Ev,v,w,d_Ew)&=\frac1{720}-\frac v{280},\\
 \beta_4(v,d_Ev,w,d_Ew)&=\frac1{720}+\frac{11v}{630}.
 \end{aligned}
\end{equation}
There are five planar binary trees. Grouping them by the three
root splittings, their projected contributions for the first input
are $0$, $1/720-v/280$, $0$. For the second input they are
$1/360+v/72$, $-1/720+v/280$, $0$.
Each is evaluated by the same radial moment, and their sums give
\eqref{mht:eq:finite-quartic-products}.

\section{Algebraic recovery and analytic convergence}

The choice of symplectic exhaustion gives exact recovery on finite
algebraic expressions. It need not control the norms of its projections.

\begin{proposition}\label{mht:prop:finite-recovery}
Fix finitely many inputs in $E_{\mathrm{alg}}$ and finitely many
expressions formed from their products and $h$. Their projections
to $E_K$ equal the original expressions for every sufficiently
large $K$.

For fixed $K$ and integer $s\geq2$, give functions the $H^{s+1}$
norm and one-forms the $H^s$ norm. The tree series in
\eqref{mht:eq:finite-tree-recursion} converges on a sufficiently
small field in this graded Sobolev space. Its radius can depend on $K$.
\end{proposition}

\begin{proof}
Lemma~\ref{mht:lem:finite-symplectic} keeps every intermediate
coefficient in $E_{\mathrm{alg}}$. A sufficiently large $V_K$
contains the acyclic components of this finite collection. Projection
fixes them, and the complementary homotopy vanishes on them. This
proves eventual recovery, including vanishing of higher corrections
on the fixed collection when every required intermediate coefficient
has entered the finite complex.

The Green operator gains one derivative. The identity
$\tau_E(\eta)=\kappa j(G\eta)$ makes the trace continuous in these
norms, since first-jet evaluation is continuous on $H^{s+1}$ for
$s>1$ in real dimension two. Multiplication is continuous, and the
finite-rank projections and $H_K$ are bounded.

Let $C$ bound
multiplication and let $C_{n-1}$ be the number of planar binary
trees with $n$ inputs. Then
\[
 \|I_n\|\leq C_{n-1}C^{n-1}\|H_K\|^{n-1}\|i_K\|^n,
 \qquad C_{n-1}\leq4^{n-1}.
\]
For $\beta_n$, replace the last homotopy by $\mathsf p_K$. Both sums are
bounded by a geometric series for a sufficiently small input.
\end{proof}

No bound uniform in $K$ has been used. The proposition therefore
separates eventual algebraic recovery from analytic removal of the
cutoff. A family with quantitative projection estimates is needed
for the latter.

\section{Backgrounds and artificial singularities}

Tensor $V$ and $W$ with $\mathfrak g=\mathfrak{gl}_N(\mathbb C)$.
Extend $p_0,p_1,h$ entrywise, writing $P_0,P_1$ for the matrix
projections. For a residual matrix $J$, set
\[
 d_J=d_E+\eta_0\operatorname{ad}_J,
 \qquad M_J=hd_J|_V=1+h\eta_0\operatorname{ad}_J|_V.
\]
The full contraction of
Theorem~\ref{mht:thm:background-contraction} gives
$M_J^{-1}=h_Jd_E|_V$ at every finite $J$. A fixed compression of
this operator need not be invertible.

\begin{example}\label{mht:ex:compressed-pole}
In a root direction with $\operatorname{ad}_J=\lambda$, compression
to the modes \eqref{mht:eq:finite-two-modes} gives
\begin{equation}\label{mht:eq:compressed-pole}
 M_{J,1}=1+\lambda\Omega^{-1}S
       =\begin{pmatrix}1&-\lambda\\\lambda/2&1\end{pmatrix},
 \qquad \det M_{J,1}=1+\lambda^2/2.
\end{equation}
The compression is singular at $\lambda=i\sqrt2$. For
$J=j\operatorname{diag}(1,-1)$, the two opposite root spaces
give determinant $(1+2j^2)^2$, and the two Cartan directions
contribute determinant one. The matrix trace pairs opposite roots.
One scalar root block is not itself a skew kinetic form.
\end{example}

The mode $w$ in this example mixes rotation weights $1$ and $-1$.
The singularity depends on the chosen subspace. To avoid it for an
arbitrary finite subspace, first change the full kinetic form to a
constant one and then transport that subspace.

\section{An entire family of cyclic frames}

Fix an integer $s\geq2$ and a smooth metric on $C$. On the scalar
Sobolev space
\[
 V_s=\{f\in H^{s+1}(C):j(f)=0,\ p_0(f)=0\}
\]
define $Xf=h(\eta_0f)$.

\begin{lemma}\label{mht:lem:quasinilpotent-green}
The operator $X:V_s\to V_s$ is compact and has spectrum $\{0\}$.
For every finite matrix $J$ and
$\alpha\in\{-1,-\tfrac12,\tfrac12\}$, the series
\begin{equation}\label{mht:eq:entire-binomial-frame}
 (1+X\otimes\operatorname{ad}_J)^\alpha
   =\sum_{n\geq0}\binom\alpha n X^n\otimes(\operatorname{ad}_J)^n
\end{equation}
converges in operator norm, locally uniformly in the entries of $J$.
It is an entire operator-valued function of those entries.
\end{lemma}

\begin{proof}
Multiplication by $\eta_0$ is smooth and $h$ gains one derivative.
Its finite-rank corrections obey the trace estimate in
Proposition~\ref{mht:prop:finite-recovery}. Hence $X$ maps $V_s$
continuously into $H^{s+2}$, and Rellich compactness proves that
it is compact on $V_s$.

Suppose $Xv=\mu v$ with $\mu\ne0$. The gain of regularity in this
equation makes $v$ smooth. Put $\lambda=-\mu^{-1}$. Since
$hd_Ev=v$, one has $h(d_E+\lambda\eta_0)v=0$. Moreover
\[
 p_1((d_E+\lambda\eta_0)v)=\lambda p_0(v)=0.
\]

The form lies in $W$, where $h$ is injective, so
$(d_E+\lambda\eta_0)v=0$. Multiplication by $e^{\lambda f_0}$
gives a holomorphic function on $\mathbb P^1$, hence
$v=ce^{-\lambda f_0}$. The first jet is $j(v)=-\lambda c$.
The Neumann condition forces $c=0$, excluding every nonzero
eigenvalue. Compactness and infinite dimension imply
$\operatorname{spec}(X)=\{0\}$.

The spectral-radius formula gives
$\lim_n\|X^n\|^{1/n}=0$. On each compact set of matrices,
$\|\operatorname{ad}_J\|$ is bounded. The binomial coefficients
have at most polynomial growth, so the root test gives normal
convergence of \eqref{mht:eq:entire-binomial-frame}. Each summand
is polynomial in $J$, proving the entire dependence.
\end{proof}

On the matrix-valued acyclic space put
\[
 \omega_J(v,w)=\tau_E\operatorname{tr}(v d_Jw),\qquad
 \omega_0(v,w)=\tau_E\operatorname{tr}(v d_Ew).
\]
Both forms are skew by cyclicity, and
\begin{equation}\label{mht:eq:kinetic-self-adjoint}
 \omega_J(v,w)=\omega_0(v,M_Jw),\qquad
 \omega_0(M_Jv,w)=\omega_0(v,M_Jw).
\end{equation}
Define the full frame
\begin{equation}\label{mht:eq:full-cyclic-frame}
 \mathcal T_J=M_J^{-1/2}
 =1-\frac12A_J+\frac38A_J^2-\frac5{16}A_J^3+\cdots,
 \qquad A_J=X\otimes\operatorname{ad}_J.
\end{equation}
The scalar projection $\Pi_K$ is now tensored with the identity on
$\mathfrak g$. Set
\begin{align}
 V_K(J)&=\mathcal T_J(V_K\otimes\mathfrak g),&
 \Pi_K^J&=\mathcal T_J\Pi_K\mathcal T_J^{-1},\notag\\
 E_K^0(J)&=\mathfrak g\oplus V_K(J),&
 E_K^1(J)&=\eta_0\mathfrak g\oplus d_JV_K(J).
 \label{mht:eq:moving-finite-complex}
\end{align}
The residual summand is the full complex
$[\mathfrak g\xrightarrow{\operatorname{ad}_J}\mathfrak g\eta_0]$.

\begin{theorem}\label{mht:thm:moving-finite-contraction}
For every finite $J$,
\begin{equation}\label{mht:eq:moving-darboux}
 \omega_J(\mathcal T_Jv,\mathcal T_Jw)=\omega_0(v,w).
\end{equation}
The complex \eqref{mht:eq:moving-finite-complex} has a nondegenerate
cyclic pairing and the cyclic strong contraction
\begin{align}
 \mathsf p_{K,J}^0f&=P_0f+\Pi_K^J(f-P_0f),\notag\\
 \mathsf p_{K,J}^1\eta&=\eta_0P_1\eta+d_J\Pi_K^Jh_J\eta,&
 H_{K,J}\eta&=(1-\Pi_K^J)h_J\eta.
 \label{mht:eq:moving-finite-contraction}
\end{align}
The homotopy is zero on degree zero. These maps depend analytically
on all entries of $J$ and preserve smooth vectors. Their formal
coefficients involve finitely many Green operations and commutators,
with no inverse eigenvalue difference.
\end{theorem}

\begin{proof}
The binomial series and \eqref{mht:eq:kinetic-self-adjoint} give
$\mathcal T_J^\dagger M_J\mathcal T_J=1$, first near zero by
series multiplication and then everywhere by analyticity.
Here $\dagger$ is the adjoint for $\omega_0$. This proves
\eqref{mht:eq:moving-darboux}. Thus $\Pi_K^J$ is the
$\omega_J$-orthogonal projection onto $V_K(J)$.

The residual summand remains orthogonal to the acyclic summand.
On $W$, conjugate the degree-zero projection by the isomorphism
$d_J:V\to W$, whose inverse is $h_J$. This gives the degree-one
projection in \eqref{mht:eq:moving-finite-contraction}. The inverse
differential on the remaining orthogonal summand is $H_{K,J}$.
The decomposition proves the contraction identity, adjointness,
cyclicity and all side conditions. The residual differential is
retained because $P_1d_Jf=[J,P_0f]$.

The series converge in every $V_s$, and their limits agree under
Sobolev inclusion. Hence they preserve smooth vectors. Near zero,
$h_J=M_J^{-1}h$, so every formal coefficient uses only finitely
many copies of $h\eta_0\operatorname{ad}_J$. This also describes
the coefficients of the transported projection and homotopy.
\end{proof}

Taking a square root of the compressed matrix
\eqref{mht:eq:compressed-pole} would retain its singularities.
The full frame supplies a nondegenerate moving subspace at every
finite background, including nonsemisimple matrices.

\section{The parameter connection and its curvature}

Let $d_{\mathrm{par}}$ differentiate the entries of $J$, and put
$\mathcal R_J=\mathcal T_J^{-1}=M_J^{1/2}$. The frame connection is
\begin{equation}\label{mht:eq:frame-parameter-connection}
 \nabla=d_{\mathrm{par}}+\mathcal A,
 \qquad \mathcal A=\mathcal R_J^{-1}d_{\mathrm{par}}\mathcal R_J.
\end{equation}
This differential acts on parameters, separately from $d_J$ on
curve coefficients.

\begin{proposition}\label{mht:prop:frame-connection}
The connection \eqref{mht:eq:frame-parameter-connection} is flat
and compatible with the varying form:
\[
 d_{\mathrm{par}}\mathcal A+\mathcal A\wedge\mathcal A=0,
 \qquad d_{\mathrm{par}}\omega_J
       =\mathcal A^{\mathsf T}\omega_J+\omega_J\mathcal A.
\]
In contrast, for $\theta=M_J^{-1}d_{\mathrm{par}}M_J$,
\begin{equation}\label{mht:eq:half-connection-curvature}
 F_{\theta/2}=-\frac14\theta\wedge\theta,
 \qquad
 \mathcal A=\frac12\theta+
       \frac18[A_J,d_{\mathrm{par}}A_J]
       +O(A_J^2d_{\mathrm{par}}A_J).
\end{equation}
The last remainder has at least two undifferentiated factors $A_J$.
\end{proposition}

\begin{proof}
The Maurer--Cartan identity for $\mathcal R_J^{-1}d_{\mathrm{par}}\mathcal R_J$
gives flatness. Differentiating
$\omega_J=\mathcal R_J^{\mathsf T}\omega_0\mathcal R_J$ proves
compatibility. The same identity for $\theta$ gives the curvature
of $\theta/2$. Expanding $\mathcal R_J=1+A_J/2-A_J^2/8+\cdots$
and its inverse gives, through the displayed order,
\[
 \mathcal A=\frac12d_{\mathrm{par}}A_J
       -\frac38A_Jd_{\mathrm{par}}A_J
       -\frac18(d_{\mathrm{par}}A_J)A_J+\cdots.
\]
Subtract $\theta/2=(d_{\mathrm{par}}A_J-A_Jd_{\mathrm{par}}A_J)/2+\cdots$
to obtain the commutator term.
\end{proof}

The curvature is nonzero for the curve operator itself. On the
mode $v$ of \eqref{mht:eq:finite-two-modes}, Green integration gives
\begin{equation}\label{mht:eq:native-frame-curvature}
 X^2v=\frac{z^2(1+3z\bar z)}{12(1+z\bar z)^3},\qquad
 F_{\theta/2}|_{J=0}(U,V)
       =-\frac14X^2\otimes\operatorname{ad}_{[U,V]}.
\end{equation}
For $U=E_{12}$ and $V=E_{21}$, the commutator is
$\operatorname{diag}(1,-1)$ and its adjoint action is nonzero.
Thus the half-logarithmic formula already fails to give a flat
connection in the native matrix family.

\section{Source antifields and change of density}

A parameter-dependent change of finite field coordinates also
changes the antifields of those parameters. Let $q$ be graded
coordinates, let $J^i$ be even parameters, and give their odd
cotangent coordinates the degrees
$|q^+|=-1-|q|$ and $|J_i^+|=-1$.

Let $T(J)$ be an invertible degree-preserving formal matrix,
with $T(0)=1$, and set $q=T(J)Q$. Its cotangent lift is determined
by the one-form identity
\begin{equation}\label{mht:eq:moving-source-cotangent}
 q^+\delta q+J_i^+\delta J^i
       =Q^+\delta Q+\widetilde J_i^+\delta J^i.
\end{equation}
Dual coordinates are placed on the left. Expanding the left side gives
\begin{equation}\label{mht:eq:moving-antifield-shift}
 Q^+=q^+T,\qquad
 \widetilde J_i^+=J_i^++q^+(\partial_iT)Q.
\end{equation}
For the moving cutoff, $T$ is a finite coordinate comparison between
chosen transported bases. No infinite-dimensional determinant of
$\mathcal T_J$ is required.

\begin{proposition}\label{mht:prop:moving-semidensity}
Write $b=\operatorname{Ber}T$, with formal branch $b(0)=1$.
The odd cotangent lift has Berezinian $b^2$ and sends a coordinate
semidensity to $b$ times that semidensity. If a master weight is
$e^{I/\hbar}$ times the old coordinate semidensity, its new action is
\begin{equation}\label{mht:eq:moving-density-action}
 I'=I\circ F+\hbar\log\operatorname{Ber}T,
\end{equation}
where $F$ is the full coordinate change including
\eqref{mht:eq:moving-antifield-shift}.
\end{proposition}

\begin{proof}
The base Jacobian in $(Q,J)$ is triangular and has Berezinian $b$.
The inverse transpose on the dual variables would have reciprocal
Berezinian, but parity reversal for the odd cotangent space inverts
it again. The full Berezinian is therefore $b^2$. Its normalized
square root is $b$, and absorbing that factor into the exponential
gives \eqref{mht:eq:moving-density-action}.
\end{proof}

For one even coordinate, $q=e^JQ$ gives $q^+=e^{-J}Q^+$ and
$\widetilde J^+=J^++Q^+Q$. The coordinate Berezinian is $e^{2J}$,
and the action acquires $\hbar J$. If equal even and odd copies
carry the same matrix $T$, their base Berezinians cancel. Their
source-antifield shifts still remain. A varying cutoff therefore
requires its cotangent lift even when its density contribution is one.
